% GENERATED FILE. Edit canonical lessons, then run npm run generate:books.
% canonical-source-hash: fnv1a64:1581d4d79bd1383f
% canonical-lessons: GE-C387-kontostand, GE-C387-guthaben, GE-C387-ueberziehen, GE-C387-abbuchen, GE-C387-betrag

\chapter{Auf dem Konto}
\label{ch:ge-auf-dem-konto}
\hlchaptermodality{387}

\begin{chapteropening}
\textbf{By the end of this chapter:} \emph{I can talk about my bank balance and credit, going overdrawn, money taken from my account, and an amount due.}

Complete the last lesson of chapter 387: I can talk about my bank balance and credit, going overdrawn, money taken from my account, and an amount due.
\end{chapteropening}

\section[der Kontostand]{der Kontostand — the bank balance}

\label{lesson:GE-C387-kontostand}



\begin{quote}
Before the new one: say the German for to be liable, then the German for liability.
\end{quote}

\subsection*{You'll want to know: der Kontostand}
\textbf{der Kontostand} — \textquotedblleft{}the bank balance\textquotedblright{}.

\begin{quote}
Jeden Morgen schaut sie auf ihren Kontostand.
\end{quote}

\begin{quote}
\emph{Every morning she checks her bank balance.}
\end{quote}

Plural: \emph{die Kontostände}.

\begin{cousinweb}
\emph{Kontostand} joins \emph{Konto} \textquotedblleft{}account\textquotedblright{}, borrowed from Italian \emph{conto} \textquotedblleft{}reckoning\textquotedblright{} (EtymDB), to \emph{Stand} \textquotedblleft{}state, standing\textquotedblright{}, from the family of \emph{stehen} \textquotedblleft{}to stand\textquotedblright{}: where your account stands.
\end{cousinweb}

\begin{grammarlens}[title={What you've built}]
One more word for telling, explaining and describing at B1.
\end{grammarlens}

\subsection*{Your turn}
\begin{itemize}
\raggedright
  \item \emph{Say it:} \emph{der Kontostand}
  \item \emph{Say it:} \emph{der Kontostand}, once more
  \item \emph{Say it:} a sentence of your own with \emph{der Kontostand}
  \item \emph{Recall:} say the German for to be liable, then the German for liability, then say \emph{der Kontostand} again
\end{itemize}

\begin{tcolorbox}[breakable,colback=teal!4,colframe=teal!35!black,title={Before you move on}]
What does der Kontostand mean? (\textquotedblleft{}The bank balance\textquotedblright{}.) Say it once more.
\end{tcolorbox}
\section[das Guthaben]{das Guthaben — the credit (balance)}

\label{lesson:GE-C387-guthaben}



\begin{quote}
Before the new one: say the German for liability, then the German for the bank balance.
\end{quote}

\subsection*{You'll want to know: das Guthaben}
\textbf{das Guthaben} — \textquotedblleft{}the credit (balance)\textquotedblright{}.

\begin{quote}
Auf meinem Handy habe ich noch zehn Euro Guthaben.
\end{quote}

\begin{quote}
\emph{I still have ten euros of credit on my mobile.}
\end{quote}

Plural: \emph{die Guthaben}. Literally what you have to your good.

\begin{cousinweb}
\emph{Guthaben} joins \emph{gut} \textquotedblleft{}good\textquotedblright{}, Old High German \emph{guot} (EtymDB), to \emph{haben} \textquotedblleft{}to have\textquotedblright{}: the sum that stands to your good.
\end{cousinweb}

\begin{grammarlens}[title={What you've built}]
One more word for telling, explaining and describing at B1.
\end{grammarlens}

\subsection*{Your turn}
\begin{itemize}
\raggedright
  \item \emph{Say it:} \emph{das Guthaben}
  \item \emph{Say it:} \emph{das Guthaben}, once more
  \item \emph{Say it:} a sentence of your own with \emph{das Guthaben}
  \item \emph{Recall:} say the German for liability, then the German for the bank balance, then say \emph{das Guthaben} again
\end{itemize}

\begin{tcolorbox}[breakable,colback=teal!4,colframe=teal!35!black,title={Before you move on}]
What does das Guthaben mean? (\textquotedblleft{}The credit (balance)\textquotedblright{}.) Say it once more.
\end{tcolorbox}
\section[überziehen]{überziehen — to overdraw (an account); to overrun}

\label{lesson:GE-C387-ueberziehen}



\begin{quote}
Before the new one: say the German for the bank balance, then the German for the credit.
\end{quote}

\subsection*{You'll want to know: überziehen}
\textbf{überziehen} — \textquotedblleft{}to overdraw (an account); to overrun\textquotedblright{}.

\begin{quote}
Er hat sein Konto um zweihundert Euro überzogen.
\end{quote}

\begin{quote}
\emph{He has overdrawn his account by two hundred euros.}
\end{quote}

In this sense inseparable: \emph{er überzieht, er hat überzogen}. Separable when you pull something on over: \emph{Er zieht eine Jacke über.}

\begin{cousinweb}
\emph{überziehen} is built from \emph{über} \textquotedblleft{}over\textquotedblright{} and \emph{ziehen} \textquotedblleft{}to pull, to draw\textquotedblright{}, inherited from Old High German \emph{ziohan} and Proto-Germanic \emph{teuhaną} (EtymDB): you draw out more than is there.
\end{cousinweb}

\begin{grammarlens}[title={What you've built}]
One more word for telling, explaining and describing at B1.
\end{grammarlens}

\subsection*{Your turn}
\begin{itemize}
\raggedright
  \item \emph{Say it:} \emph{überziehen}
  \item \emph{Say it:} \emph{überziehen}, once more
  \item \emph{Say it:} a sentence of your own with \emph{überziehen}
  \item \emph{Recall:} say the German for the bank balance, then the German for the credit, then say \emph{überziehen} again
\end{itemize}

\begin{tcolorbox}[breakable,colback=teal!4,colframe=teal!35!black,title={Before you move on}]
What does überziehen mean? (\textquotedblleft{}To overdraw (an account); to overrun\textquotedblright{}.) Say it once more.
\end{tcolorbox}
\section[abbuchen]{abbuchen — to debit, to take (from an account)}

\label{lesson:GE-C387-abbuchen}



\begin{quote}
Before the new one: say the German for the credit, then the German for to overdraw.
\end{quote}

\subsection*{You'll want to know: abbuchen}
\textbf{abbuchen} — \textquotedblleft{}to debit, to take (from an account)\textquotedblright{}.

\begin{quote}
Die Miete wird jeden Monat automatisch abgebucht.
\end{quote}

\begin{quote}
\emph{The rent is taken from the account automatically every month.}
\end{quote}

Separable: \emph{die Bank hat das Geld abgebucht}. \emph{Die Abbuchung} — the debit.

\begin{cousinweb}
\emph{abbuchen} is built from \emph{ab} \textquotedblleft{}off\textquotedblright{} and \emph{buchen} \textquotedblleft{}to enter in a book\textquotedblright{}, made from \emph{Buch} \textquotedblleft{}book\textquotedblright{}, Old High German \emph{buoh} and Proto-Germanic \emph{bōks} (EtymDB).
\end{cousinweb}

\begin{grammarlens}[title={What you've built}]
One more word for telling, explaining and describing at B1.
\end{grammarlens}

\subsection*{Your turn}
\begin{itemize}
\raggedright
  \item \emph{Say it:} \emph{abbuchen}
  \item \emph{Say it:} \emph{abbuchen}, once more
  \item \emph{Say it:} a sentence of your own with \emph{abbuchen}
  \item \emph{Recall:} say the German for the credit, then the German for to overdraw, then say \emph{abbuchen} again
\end{itemize}

\begin{tcolorbox}[breakable,colback=teal!4,colframe=teal!35!black,title={Before you move on}]
What does abbuchen mean? (\textquotedblleft{}To debit, to take (from an account)\textquotedblright{}.) Say it once more.
\end{tcolorbox}
\section[der Betrag]{der Betrag — the amount, the sum}

\label{lesson:GE-C387-betrag}



\begin{quote}
Before the new one: say the German for to overdraw, then the German for to debit.
\end{quote}

\subsection*{You'll want to know: der Betrag}
\textbf{der Betrag} — \textquotedblleft{}the amount, the sum\textquotedblright{}.

\begin{quote}
Der Betrag ist bis Ende des Monats fällig.
\end{quote}

\begin{quote}
\emph{The amount is due by the end of the month.}
\end{quote}

Plural: \emph{die Beträge}. \emph{Der Gesamtbetrag} — the total.

\begin{cousinweb}
\emph{Betrag} is the noun of \emph{betragen} \textquotedblleft{}to amount to\textquotedblright{}, built with \emph{be-} on \emph{tragen} \textquotedblleft{}to carry\textquotedblright{}, Proto-Germanic \emph{draganą} (EtymDB): what a bill carries.
\end{cousinweb}

\begin{grammarlens}[title={What you've built}]
One more word for telling, explaining and describing at B1.
\end{grammarlens}

\subsection*{Your turn}
\begin{itemize}
\raggedright
  \item \emph{Say it:} \emph{der Betrag}
  \item \emph{Say it:} \emph{der Betrag}, once more
  \item \emph{Say it:} a sentence of your own with \emph{der Betrag}
  \item \emph{Recall:} say the German for to overdraw, then the German for to debit, then say \emph{der Betrag} again
\end{itemize}

\begin{tcolorbox}[breakable,colback=teal!4,colframe=teal!35!black,title={Before you move on}]
What does der Betrag mean? (\textquotedblleft{}The amount, the sum\textquotedblright{}.) Say it once more.
\end{tcolorbox}
